\documentclass[aps,prx,preprint,superscriptaddress,nofootinbib,longbibliography]{revtex4-2}
\usepackage[T1]{fontenc}
\usepackage{amsmath,amssymb,mathtools,booktabs,microtype,hyperref}
\newcommand{\E}{\mathbb E}\newcommand{\Prob}{\mathbb P}
\begin{document}
\title{Supplemental Material for ``Throughput Is Not Enough: Synchronization and Loss Constraints for GKP-State Factories in High-Rate Photonic Repeaters''}
\author{Ayush Nadiger}
\affiliation{Department of Electrical and Computer Engineering, University of Massachusetts Amherst, Amherst, Massachusetts 01003, USA}
\date{August 31, 2026}
\maketitle

\section{Inverse-design convention}
The numerical benchmark in the main text is deliberately separated from the resource-success labels reported by Shiina \emph{et al.} Their $>0.999$ and $>0.9999$ initial-resource counts already refer to simultaneous preparation success across the repeater chain and are therefore not reinterpreted as an independent downstream conditional probability.

The standalone factory benchmark instead declares
\begin{equation}
N=3800,\qquad S=111,\qquad p_F=0.999,
\end{equation}
where $N$ is a rounded payload-scale anchor motivated by the approximately 3800 initial GKP qubits per station per protocol run reported at a representative high-rate repeater point, $S$ is a representative synchronization scale, and $p_F$ is an independently chosen factory-readiness target. No present finite-energy GKP source is assumed to satisfy this contract.

A generic reliability composition rule remains valid when two events are genuinely distinct: if a factory-qualified event $F$ is followed by a downstream event $C$ with $\Prob(C\mid F)=p_C$, then $\Prob(F)\ge p_{\rm tot}/p_C$ is sufficient for total readiness $p_{\rm tot}$. We do not use that decomposition for the benchmark above.

\section{Independent synchronized service}
For iid independent station counts $A_s$, synchronized factory readiness requires
\begin{equation}
\Prob(A_1\ge N,\ldots,A_S\ge N)\ge p_F,
\end{equation}
so the local target is
\begin{equation}
p_{\rm loc}=p_F^{1/S}=0.9999909865.
\end{equation}
The associated one-station shortage probability is approximately $9.0135\times10^{-6}$ per cycle.

If $M$ independent source opportunities each yield one service-qualified state with probability $y$, then $A\sim\mathrm{Binomial}(M,y)$ and the minimum $M$ is obtained by exact integer tail inversion. In the low-yield/high-parallelism limit, $A\sim\mathrm{Poisson}(\lambda)$ with $\lambda=My$ and the requirement becomes
\begin{equation}
Q(N,\lambda)=\Prob[A\ge N]\ge p_{\rm loc}.
\end{equation}
Monotone inversion of the exact Poisson survival function at $N=3800$ gives
\begin{equation}
\boxed{\lambda_*=4070.1489},
\end{equation}
with count overhead
\begin{equation}
\frac{\lambda_*-N}{N}=0.0710918=7.10918\%.
\end{equation}
At an illustrative factory clock of $10^5\,\mathrm{s}^{-1}$ this corresponds to $4.07015\times10^8$ accepted states/s/station. The clock is a scenario variable; the per-cycle tail law does not require it.

\section{Gamma--Poisson common-mode benchmark}
Let one station-cycle have a random multiplicative rate factor $Z$ with $\E Z=1$ and
\begin{equation}
A\mid Z=z\sim\mathrm{Poisson}(\lambda z).
\end{equation}
For a Gamma rate factor with shape $k$ and scale $1/k$,
\begin{equation}
\operatorname{Var}(Z)=1/k.
\end{equation}
An RMS fractional rate fluctuation $r$ therefore corresponds to $k=r^{-2}$. Poisson--Gamma mixing gives
\begin{equation}
\Prob(A=a)=\frac{\Gamma(a+k)}{\Gamma(k)a!}
\left(\frac{k}{k+\lambda}\right)^k
\left(\frac{\lambda}{k+\lambda}\right)^a,
\end{equation}
with mean $\lambda$ and variance $\lambda+\lambda^2/k$.

Exact lower-tail inversion at the same $p_{\rm loc}$ gives
\begin{align}
\lambda_{5\%}&=4798.8595, & \frac{\lambda_{5\%}-N}{N}&=26.2858\%,\\
\lambda_{10\%}&=6083.0401, & \frac{\lambda_{10\%}-N}{N}&=60.0800\%.
\end{align}
These are distribution-specific benchmarks, not universal jitter penalties. Their purpose is to demonstrate that moderate common-mode variation can dominate independent shot noise when the service event probes an extreme lower tail.

If only the mean and variance of accepted count are known, Cantelli's inequality gives
\begin{equation}
\Prob(A<N)\le\frac{\sigma_A^2}{\sigma_A^2+(\mu_A-N)^2},
\end{equation}
which is far too loose to certify the required local shortage probability at modest overhead. RMS stability alone is therefore insufficient; a physical factory must provide either a measured lower tail/joint distribution or a justified stochastic model.

\section{Pooling versus labeled near-determinization}
Suppose $N$ labeled downstream ports are each made ready independently with probability $1-\delta$. Requiring every port to be ready with probability at least $1-\epsilon$ gives
\begin{equation}
(1-\delta)^N\ge1-\epsilon,
\end{equation}
and therefore $\delta\lesssim\epsilon/N$ in the small-failure regime. In common independent-retry constructions, driving a labeled failure probability to $O(1/N)$ requires logarithmically growing retry/resource depth. A pooled inventory instead pays one aggregate count-tail constraint. The statistical advantage must be charged against the physical switch depth, insertion loss, and control bandwidth required to realize that fungibility.

\section{One-cycle FIFO reserve and quasi-stationary benchmark}
Let $X_t\sim\mathrm{Poisson}(\lambda)$ be new accepted production, $B_t\in\{0,\ldots,B_{\max}\}$ the reserve at cycle start, and demand $N$. A cycle is served when $B_t+X_t\ge N$, after which
\begin{equation}
B_{t+1}=\min\{B_{\max},B_t+X_t-N\}.
\end{equation}
Failure makes the nonfailure transition matrix substochastic:
\begin{equation}
Q_{ij}=\sum_{x=0}^{\infty}e^{-\lambda}\frac{\lambda^x}{x!}
\mathbf 1\!\left[j=\min\{B_{\max},i+x-N\}\right]
\mathbf 1[i+x\ge N].
\end{equation}

For $B_{\max}=N$, let $q$ be the normalized left quasi-stationary distribution of $Q$ and $\rho$ its dominant eigenvalue. Starting in that survivor-conditioned operating distribution, one station survives $H$ more cycles with probability $\rho^H$. We make the benchmark horizon explicit and impose
\begin{equation}
\rho^{10S}\ge0.999
\end{equation}
for ten consecutive synchronized cycles across $S=111$ stations.

Monotone inversion yields
\begin{equation}
\boxed{\lambda_{\rm buf}\simeq3804.633},
\end{equation}
corresponding to
\begin{equation}
\frac{\lambda_{\rm buf}-N}{N}\simeq0.001219=0.1219\%.
\end{equation}
This number is not a stationary no-failure guarantee over an infinite horizon; it is the declared ten-cycle quasi-stationary service benchmark.

Under FIFO, old reserve states are consumed first. If $U_t=\min\{B_t,N\}$ is the number of one-cycle-old states consumed, then
\begin{equation}
f_{\rm old}=\frac{\E_q U_t}{N}\simeq0.9015.
\end{equation}
The same quasi-stationary calculation gives mean overflow discarded after a served cycle of approximately 4.64 states. Thus the count-only benefit is purchased by making most consumed resources one cycle old.

If age changes the GKP quality variable or analog side information, the count service event must be replaced by a quality-qualified event such as
\begin{equation}
\#\{i:Q_i(A_i)\ge Q_{\min},\ \eta_i\ge\eta_{\min}\}\ge N.
\end{equation}
The present count model deliberately leaves that physical quality map to the source/refinery and repeater decoder.

\section{Routing-loss inversion}
If each accepted state traverses $d$ statistically similar routing stages with power transmission $\eta_{\rm sw}$,
\begin{equation}
\eta_{\rm route}=\eta_{\rm sw}^d.
\end{equation}
In decibels, $L_{\rm route}=dL_{\rm sw}$, so any repeater-level routing-loss allowance implies
\begin{equation}
\boxed{L_{\rm sw}\le L_{\rm route}/d.}
\end{equation}
For representative large-port depth $d=17$, even a modest aggregate routing allowance can imply a per-stage target of only a few millidecibels. This is an inverse-design constraint, not an assertion that a present switch platform reaches it.

\section{Reproducibility protocol}
Every stochastic number in the corrected main text follows from the following auditable sequence:
\begin{enumerate}
\item declare $N=3800$, $S=111$, and $p_F=0.999$ independently of the source paper's chain-wide resource-success labels;
\item compute $p_{\rm loc}=p_F^{1/S}$ and invert the exact Poisson survival function;
\item for common-mode benchmarks set Gamma shape $k=r^{-2}$ and invert the exact negative-binomial survival function;
\item for buffering construct the substochastic reserve kernel with $B_{\max}=N$, compute its dominant eigenvalue/vector, and solve $\rho^{10S}=0.999$;
\item evaluate old-state consumption and overflow in the same quasi-stationary ensemble;
\item convert any downstream routing-loss allowance into a per-stage target by division of the decibel budget by routing depth.
\end{enumerate}
A physical follow-up then supplies the joint distribution $P(N_{\rm acc},Q,A,\eta_{\rm route})$ and composes it with the actual repeater construction and analog decoder.

\begin{thebibliography}{10}
\bibitem{Shiina2026} R. Shiina, K. Goodenough, N. Arnold, and F. Rozp\k{e}dek, High-Rate and Resource-Efficient All-Photonic Quantum Repeater Architectures with 9 km Repeater Spacing, arXiv:2606.25314v2 (2026).
\bibitem{Larsen2025} M. V. Larsen \emph{et al.}, Integrated photonic source of Gottesman--Kitaev--Preskill qubits, \emph{Nature} \textbf{642}, 587--591 (2025).
\end{thebibliography}
\end{document}
